\section{Small sizes}\label{sec:small}

This section is numerical, except for \cref{prop:aperiodic}, which is a formal statement about one explicit code, and
for the results of \cref{sec:two-class} and of the subsequent memo~\cite{ClassUniformMemo2026} that are cited at the
end of \cref{rem:small} and at the end of the section. It reports exact computations over $\F_2$ for class-uniform rules
(\cref{def:class-uniform}) on a cycle. No proof in this paper uses them.

\subsection{Method}

For a partition of $\Z/d$ and a target dimension $D$, the existence of a class-uniform rule with a code of dimension at
least $D$ is a rank condition: the $d^2\times d^2$ matrix of $\sw-K$, whose row $(c,r)$ is determined by the functional
of the class of $c$ for the coordinate~$r$, must have rank at most $d^2-D$. We decide it in two ways.
\begin{itemize}[leftmargin=1.5em]
\item \emph{Exhaustively}, for $d\le4$, by enumerating subspaces and testing the criterion that a subspace $\X$ is a code
if and only if, for every class $J$, $\sum_{c\in J}\sigma^cX_c=0$ with $X_c\in\X$ implies
$\sum_{c\in J}(\text{row }c\text{ of }X_c)=0$ (\leanfile{tools/exhaustive_pigeonhole_checks.py}).
\item \emph{By a SAT solver}, for $d=5,6$: the rank condition is written as a factorization of the matrix through
$\F_2^{d^2-D}$, with the second factor in reduced echelon form, and given to
CryptoMiniSat (\leanfile{tools/class_uniform_sat.py}). A satisfying assignment is a rule together with its code; the
script rebuilds the matrix from the rule, recomputes its kernel, and \leanfile{tools/verify_class_uniform_code.py}
checks the saved code against the definition, sharing no code with the search. An answer ``unsatisfiable'' rests on the
solver, and we say so wherever we report one; as controls, the solver confirms the bound of \cref{thm:TP} for $d=4$ and
for $\{0,1\},\{2,3\},\{4\}$ at $d=5$. A time-out says nothing: the solver also times out on instances that are known to
be satisfiable. Some searches were restricted to codes that contain the space $\Hk^*$ of matrices with
$x_{c'}=\sigma^{c'-c}x_c$ for all $c,c'$ in one class, which is itself a code of dimension $Md$ and, for interval
classes, is the space $\Hk$ of \cref{sec:symmetry}. This makes the search faster and can only lower the result. The
code of \cref{prop:aperiodic} was found without the restriction and contains $\Hk^*$.
\end{itemize}

\subsection{Results}

\Cref{tab:small} lists the largest dimension of a code for partitions of $\Z/d$, up to the affine maps
$c\mapsto uc+t$ of the index. For $d\le5$ every partition with $1<M<d$ classes is listed; for $d=6$, the partitions
into two classes and the nine partitions into three classes for which the bound exceeds $(d^2+Md)/2$. ``Bound'' is the best value of \cref{thm:TP} over shift sets
of at most four elements, rounded down. Every code counted in the last column for $d\ge5$ is saved in
\leanfile{tools/data/}. The script \leanfile{tools/small_sizes.py} prints the table from these files, with the bounds
recomputed and every saved code checked again, and lists the partitions; \leanfile{tools/README.md} records the runs
of the solver, including those that timed out.

\begin{table}[ht]
\centering\small
\begin{tabular}{@{}cllccl@{}}
\toprule
$d$ & classes & pairing & $\lfloor(d^2+Md)/2\rfloor$ & bound & largest code \\
\midrule
3 & $\{0\},\{1,2\}$ & 6 & 7 & 7 & $=7$ \\
\midrule
4 & $\{0\},\{1,2,3\}$ & 10 & 12 & 12 & $=12$ \\
  & $\{0,1\},\{2,3\}$ & 10 & 12 & 12 & $=12$ \\
  & $\{0,2\},\{1,3\}$ & 10 & 12 & 12 & $=12$ \\
  & $\{0\},\{1\},\{2,3\}$ & 10 & 14 & 14 & $=14$ \\
  & $\{0\},\{2\},\{1,3\}$ & 10 & 14 & 14 & $=14$ \\
\midrule
5 & $\{0\},\{1,2,3,4\}$ & 15 & 17 & 17 & $=17$ \\
  & $\{0,1\},\{2,3,4\}$ & 15 & 17 & 17 & $=17$ \\
  & $\{0\},\{1,2\},\{3,4\}$ & 15 & 20 & 20 & $=20$ \\
  & $\{0\},\{1\},\{2,3,4\}$ & 15 & 20 & 20 & $19^{*}$ \\
  & $\{0\},\{1,4\},\{2,3\}$ & 15 & 20 & 21 & $20^{*}$ \\
  & $\{0\},\{1\},\{2\},\{3,4\}$ & 15 & 22 & 22 & $=22$ \\
\midrule
6 & $\{0\},\{1,\dots,5\}$ & 21 & 24 & 24 & $=24$ \\
  & $\{0,1\},\{2,\dots,5\}$ & 21 & 24 & 24 & $=24$ \\
  & $\{0,1,2\},\{3,4,5\}$ & 21 & 24 & 24 & $=24$ \\
  & $\{0,2,4\},\{1,3,5\}$ & 21 & 24 & 24 & $=24$ \\
  & $\{0,2\},\{1,3,4,5\}$ & 21 & 24 & 27 & $\ge24$ \\
  & $\{0,3\},\{1,2,4,5\}$ & 21 & 24 & 27 & $\ge23$ \\
  & $\{0,1,3\},\{2,4,5\}$ & 21 & 24 & 28 & $\ge25$ \\
  & $\{0,1\},\{2,4\},\{3,5\}$ & 21 & 27 & 30 & $\ge28$ \\
  & $\{0,1\},\{2,5\},\{3,4\}$ & 21 & 27 & 30 & $\ge28$ \\
  & $\{0,2\},\{1,4\},\{3,5\}$ & 21 & 27 & 30 & $\ge28$ \\
  & $\{0\},\{1,2,4\},\{3,5\}$ & 21 & 27 & 30 & $\ge28$ \\
  & $\{0\},\{1,5\},\{2,3,4\}$ & 21 & 27 & 30 & $\ge27$ \\
  & $\{0\},\{1,2,5\},\{3,4\}$ & 21 & 27 & 30 & $\ge27$ \\
  & $\{0\},\{1,3,4\},\{2,5\}$ & 21 & 27 & 30 & $\ge27$ \\
  & $\{0\},\{4\},\{1,2,3,5\}$ & 21 & 27 & 30 & $\ge27$ \\
  & $\{0\},\{3\},\{1,2,4,5\}$ & 21 & 27 & 30 & $\ge26$ \\
\bottomrule
\end{tabular}
\caption{Largest codes of class-uniform rules over $\F_2$. ``$=$'' marks an exact value: a code of that dimension
was found, and it equals the bound of \cref{thm:TP}. ``$\ge$'' marks the largest code found. A star marks a value that
was found and whose optimality rests on the solver (\cref{rem:small}).}\label{tab:small}
\end{table}

\begin{remark}[Numerical: what the table shows]\label{rem:small}
(a) \emph{Intervals.} For interval classes the bound $\lfloor(d^2+Md)/2\rfloor$ of \cref{thm:TP} is attained in every
case computed but one. For unequal pieces the codes that attain it are not block-symmetric, and at $d=4$ with pieces
$\{0\},\{1,2,3\}$ no copy rule attains it (\cref{rem:extremal}). The exception is $\{0\},\{1\},\{2,3,4\}$ at $d=5$: a
code of dimension $19$ exists; by \cref{lem:extremal} a code of dimension $20$ would contain $\Hk$; and the solver
reports that no code of dimension $20$ contains~$\Hk$. So the largest code has dimension $19$, on the word of the
solver, and the bound of \cref{thm:TP} is not attained for every partition into intervals.
(b) \emph{No almost-period, $d=5$.} For $\{0\},\{1,4\},\{2,3\}$ every single shift leaves four components, and the
bound of \cref{thm:TP} is $21$, from a set of three shifts. A code of dimension $20=(d^2+Md)/2$ exists, and the solver reports that
none of dimension $21$ does.
(c) \emph{No almost-period, $d=6$.} For $\{0,1,3\},\{2,4,5\}$ every nonzero shift leaves four components, and the best
bound of \cref{thm:TP} is $28$, from a set of three shifts. A code of dimension $25$ exists, one more than
$(d^2+Md)/2=24$
(\cref{prop:aperiodic}); searches for dimension $26$ timed out. Four partitions into three classes have codes of
dimension $28$, one more than $(d^2+Md)/2=27$.
(d) The gain over $(d^2+Md)/2$ in (c) is one dimension in each case. The data, at $d\le6$, do not distinguish a gain
of bounded size from one that grows with~$d$. For every fixed number $M\ge2$ of classes it can grow, and this is proved, not
numerical, though for other partitions than those of the table. For two classes the gain is at least $d/4-1$ at every
$d=4k\ge8$, over every field (\cref{prop:two-class}, which extends \cite[Theorem~4.8]{ClassUniformMemo2026} to
characteristic~$2$); the codes over $\F_2$ at $d=8$, $12$ and $16$ are saved with those of \cref{tab:small}
(\cref{rem:two-class}). For every $M\ge2$ there are partitions into $M$ classes with a gain of order $Md/8$
(\cref{cor:every-M}). And when the number of classes grows in proportion to $d$, the blow-ups of the code of
\cref{prop:aperiodic} are codes of dimension $25k^2$ for $2k$ classes on $\Z/6k$, a gain
of~$k^2$~\cite[Theorem~4.2 and Corollary~4.4]{ClassUniformMemo2026}.
\end{remark}

\begin{proposition}[Two classes with no almost-period beat every interval rule]\label{prop:aperiodic}
\formal{AperiodicExample.lean}{exists_classUniformCode_gt,isClassUniformCode_aperiodicCode,finrank_aperiodicCode,aperiodicCode,aperiodicCls_eq_true_iff}
For $d=6$ over $\F_2$ and the classes $\{0,1,3\}$ and $\{2,4,5\}$ of $\Z/6$ there is a code of a class-uniform rule of
dimension $25$. In particular the bound $(d^2+Md)/2$ of \cref{thm:TP} for $M$ interval classes, and of \cref{prop:copy}
for copy rules with $M$ classes of any shape, does not hold for coded rules on arbitrary classes.
\end{proposition}

\begin{outline}
The code is the span of $25$ explicit matrices, and the two decoders are explicit
(\leanfile{AperiodicExampleData.lean}, generated from a certificate by \leanfile{tools/gen_aperiodic_example.py}). That
every generator satisfies the rule at every row, $900$ equations over $\F_2$, is checked by the kernel
(\lean{aperiodicGen_row}); by linearity the span is a code (\lean{isClassUniformCode_aperiodicCode}). Each generator has
an entry equal to~$1$ where all the others vanish (\lean{aperiodicGen_pivot}), so the generators are independent and
the dimension is~$25$ (\lean{finrank_aperiodicCode}). Then $2\cdot25=50>48=6^2+2\cdot6$
(\lean{exists_classUniformCode_gt}). The certificate is used only as data.
\end{outline}

The code of \cref{prop:aperiodic} was found by the solver. It contains the $12$-dimensional space $\Hk^*$, it meets its translate by one column in exactly $2\cdot25-36=14$
dimensions, the least that \cref{lem:pigeonhole} allows, and it meets its translates by two and three columns in $18$
and $17$ dimensions (\leanfile{tools/analyse_code.py}). We do not have a description of the code itself, or a way to
enlarge it at a fixed number of classes, and one example at $d=6$ says nothing about
rates: the term $Md/2$ is a quarter of $(d^2+Md)/2$ at this size. What the example does show is that the union bound in
\cref{thm:TP} is not an artefact, and that masks without an almost-period carry a coding gain that no copy rule and no
interval rule has. \Cref{prob:aperiodic} asks how large it can be.

The subsequent memo~\cite{ClassUniformMemo2026} proves two further facts about this example, in Lean, and it finds
larger gains of a different kind.
\begin{itemize}[leftmargin=1.5em]
\item \emph{The gain does not come from characteristic~$2$.} The same two classes carry a code of dimension $25$ over
every field in which $3\ne0$, in particular over $\Q$~\cite[Theorem~3.2]{ClassUniformMemo2026}. Whether they do in
characteristic~$3$ is open.
\item \emph{The example can be blown up.} If a colouring of $\Z/d$ with $M$ classes carries a code of dimension $D$,
then a colouring of $\Z/dk$ with $kM$ classes carries a code of dimension
$k^2D$~\cite[Theorem~4.2]{ClassUniformMemo2026}. Blowing up multiplies the gain over $(d^2+Md)/2$ and the product $dM$
by the same factor $k^2$, so blowing up this one example gives growing gains only together with a growing number of
classes. Blowing up a family by a fixed factor is another matter (\cref{cor:every-M}).
\item \emph{Two classes can gain more than one dimension.} For $d=4k\ge8$ the classes $\{0,2,\dots,d-4\}$ and the rest
carry an explicit rule, each of whose decoders reads one anti-diagonal, with a code of dimension
$(d^2+2d)/2+d/4-1$. The memo proves this over fields in which $2\ne0$, in its text, and proves in Lean at $d=8$, $12$
and~$16$ that the classes carry a code of this dimension~\cite[Theorems~4.8 and~4.9]{ClassUniformMemo2026};
\cref{prop:two-class} proves it over every field, and \cref{cor:every-M} derives from it an unbounded gain for every
fixed number of classes. Unlike the code of \cref{prop:aperiodic}, these codes have a description, and it is one that
can be analysed.
\end{itemize}
